\chapter{Solver-, Ergebnis- und Zeitverträge}
\label{api:solvers}\index{Solver}\index{Ergebnis}\index{Backend}
\begingroup\raggedright
\section{Stationäre direkte APIs}
Die exakten Schlüsselwortargumente und Rückgabetypen stehen im Inventar.
\texttt{solve\_heat}, \texttt{solve\_elasticity},
\texttt{solve\_electrostatics}, \texttt{solve\_piezoelectric} nehmen Mesh,
Material(e) und feldpassende Dirichlet-Daten; optional Lasten/Elektroden/
thermische Eigenverzerrung gemäß Physik. Jeder DOF-Helfer liefert die passende
DofMap, damit inaktive Netzpunkte nicht als aktive Freiheitsgrade erscheinen.
\texttt{solve\_thermomechanics} und \texttt{solve\_thermopiezoelectric}
übertragen zunächst $T$ auf Mechanik/Piezo auf demselben Netz. Keine
reziproke Wärmeproduktion oder mechanisch-elektrische Temperaturantwort.
Thermomechanik hat keinen separaten \texttt{Model.physics}-Modus.

Sei $n$ die Anzahl aktiver Knoten, $e$ die Eingabezellzahl.
HeatResult: temperature$[n]$ K, gradient$[e,3]$ K/m, heat\_flux$[e,3]$ W/m$^2$,
reactions/residual$[n]$ W. ElasticResult: displacement$[n,3]$ m,
strain$[e,6]$ dimensionslos, stress$[e,6]$ Pa, energy$[e]$ J,
reactions/residual$[n,3]$ N, Kraftbilanz$[3]$ N, Momentbilanz$[3]$ N m.
ElectrostaticResult: potential$[n]$ V, E$[e,3]$ V/m, D$[e,3]$ C/m$^2$,
energy$[e]$ J, reactions/residual$[n]$ C. PiezoelectricResult kombiniert $u$
und $\phi$ samt feldweisen Residuen, Kraft-/Moment-/Ladungsbilanz und
Energiegrößen. Die Kopplungsergebnisse enthalten die beiden Teilresultate;
keine ungekennzeichnete gemischte Bilanz.

Tet4-Rohfelder sind zellweise konstant; Tet10-gradient/stress/E/D sind
Schwerpunktwerte mit \texttt{field\_points[e,3]} in m, keine geglätteten
Knotenfelder und keine überall konstanten quadratischen Felder. Ergebnisarrays
sind owned/readonly; Bilanzen gehören zum tatsächlich diskreten Modell.
Quasi-statische lineare kleine Deformationen; gültige Lagerung/Gauge,
positive Materialblöcke und ausreichende Geometrie sind Voraussetzungen.
Nullraum, Material-/Topologiefehler, fehlgeschlagene Residuen/Bilanzen und
Backendfehler werden diagnostiziert; kein Ergebnis mit künstlichem Erfolg.
Beispiele: \nolinkurl{examples/steady_heat.py},
\nolinkurl{examples/elastic_tension.py}, \nolinkurl{examples/parallel_plate.py},
\nolinkurl{examples/piezo_d33.py}, \nolinkurl{examples/thermopiezoelectric/run.py}.

\section{Zeitprofile und Wärmezeitintegration}
Step und Ramp beschreiben SI-Amplituden und Zeiten in s. Ramp ist außerhalb
des Intervalls konstant, innerhalb linear; endliche Grenzen und steigender
Zeitraum. TimeSeries nimmt times$[nt]$ und Werte derselben Zeilenanzahl,
explizite lineare/previous-Interpolation sowie Events. Endliche Abfragezeiten,
Ereignisseiten und keine still erfundene Extrapolation gelten laut Doc/33.
\texttt{integration\_times} verlangt streng steigende endliche Raster und
fügt erklärte Ereignisse ein; die robuste Rampen-Profilabfrage ist keine
Zusage numerisch beliebig großer Integrationszeitschritte.

\texttt{solve\_transient\_heat}: konstante positive Dichte/Wärmekapazität,
Backward Euler, Temperaturbedingungen und zeitabhängige Quellen/Flächenlasten.
TransientHeatResult bindet frames an Times in s; HeatFrame enthält aktive
Temperaturen K, Zeit und Residualdaten. HeatBalance enthält Zeitschritt s und gespeicherte Energieänderung J; es trennt
Speicherung, Quellen, vorgeschriebenen Abfluss, Konvektion und Dirichlet-Zufuhr
in W. Anfangsfeld und Sprung-/Ereigniskonventionen sind explizit; kein eigener
adaptive-Zeitintegration-/Temperaturmaterialanspruch. HeatSeriesWriter bindet
an exakt passendes Mesh/Hash und schreibt VTU/PVD/CSV/Manifest.
Minimalbeispiel: \nolinkurl{examples/transient_heat/run.py}.

\section{Piezo-Moden, Newmark und Harmonik}
\texttt{solve\_piezoelectric\_modes}: konsistente mechanische Masse,
algebraisches elektrisches Feld, statische elektrische Kondensation,
positive Moden und Massennormierung. Ergebnisfrequenzen rad/s und Hz;
modaler $u$-Vektor ist massennormiert, keine willkürliche physische Amplitude.
Keine Moden bei ungeprüften Starrkörper-/Gauge-Moden. Dichte erforderlich;
Randbedingungen homogen in der Eigenwertanalyse. Referenz ist eine begrenzte
dichte Analyse, nicht ein unbegrenzt skalierender Eigenwertdienst.

\texttt{solve\_transient\_piezoelectric}: ungedämpfte Average-Acceleration-
Newmark-DAE mit $u,v,a$ in m, m/s, m/s$^2$ und $\phi$ in V.
Felder mechanical$[nt,n,3]$, potential$[nt,n]$; Originalresiduen enthalten
Trägheit/Reaktionen. KinematicValues liefert konsistente vorgeschriebene
Ableitungen; endliche Sprünge in $u/v$ sind nicht unterstützt. Energie J,
Zeit s, feldweise Residuen und schwebende Elektrodenladung bleiben getrennt.
\texttt{solve\_harmonic\_piezoelectric}: nichtnegative einzelne Frequenz in Hz
(einschließlich DC), komplexe Spitzenwerte ohne Dämpfung; Resonanz und
ill-konditionierte Systeme werden abgelehnt. Komplexe Ergebnisse folgen
dem Inventarvertrag. Kein realer Export unter
Verlust der Imaginärteile. PiezoDynamicResult bindet die Lösung an Mesh,
Materialien und Provenienz, ModelResult an Physik/Schema.
Beispiele: \nolinkurl{examples/piezo_dynamics/run.py},
\nolinkurl{examples/tet10_modes/run.py},
\nolinkurl{examples/tet10_newmark/run.py},
\nolinkurl{examples/tet10_harmonic/run.py}; Grenzen siehe Doc/34,51,54,55.

\section{Optionale Backends und interne Reduktionsgrenze}
Der \texttt{backend}-Hook der Python-Solver ist eine experimentelle
Erweiterungsgrenze; Model/TOML bleibt SciPy. LinearBackend faktorisiert das
bereits reduzierte System, ohne selbst Lagerungen/Gauge/Elektroden zu ändern.
Feldpartition, SI-Skalierung, $(n,)$ oder $(n,nrhs)$, real/complex und
positive-definite/general müssen explizit passen. LinearSolution/SolveReport
besitzen unskalierte Werte und feldweise Residuen; konvergenzfehlerhafte
Lösungen werden nicht als Erfolg zurückgegeben. ScipyBackend, PetscBackend
und PetscFieldSplitBackend werden im Forschungs-/Backendvertrag beschrieben,
keine neue Stabilitätszusage für diese internen Adapter.

PETSc ist optional und erst bei angefordertem Backend erforderlich. Reale und
komplexe native Builds sind verschiedene Umgebungen. Zu kleines KSP-Budget,
fehlende Dependency, inkompatibler Skalartyp oder fehlgeschlagener Schur-
Preconditioner bleiben Fehler. Mathematische/physikalische Toleranzen kommen
vom jeweiligen Solver, nicht von der KSP-Erfolgsmeldung allein.
Minimalnachweise: \nolinkurl{tests/petsc/test_fieldsplit.py}
und \nolinkurl{Doc/39_PETSC_FIELDSPLIT.md}.
\endgroup
